(1) 法线向上为正。冲量动量定理在竖直方向给 \(m(v_y'-v_y)=J_n\)，且 \(v_y=-u\)、\(v_y'=eu\)，因此 \(J_n=m(1+e)u\)。该结果不依赖摩擦切向速度，假设刚性瞬时碰撞且忽略碰撞期间重力冲量。
(2) 切向冲量 \(J_t\) 使 \(v_x'=v_x+J_t/m\)，作用点位于球心下方使 \(I\omega'=I\omega+RJ_t\)，所以 \(s'=s+J_t(1/m+R^2/I)\)。要求 \(s'=0\) 得 \(J_t=-s/[1/m+R^2/I]=-2ms/7\)，进而 \(v_x'=v_x-2s/7\)、\(\omega'=\omega-5s/(7R)\)。
(3) 库仑型碰撞冲量不能超过 \(\mu J_n\)，因此恰能停止接触点切向滑动的条件为 \(2m|s|/7\le\mu m(1+e)u\)，即 \(|s|\le(7/2)\mu(1+e)u\)。超过时只能取滑移冲量 \(J_t=-\mu J_n\operatorname{sgn}(s)\)，不可继续套用上一问的无滑动式。
(4) 法向碰撞前后动能差为 \(\Delta K_n=\tfrac12m(e^2-1)u^2\)。切向包含平动和转动，其动能改变量为 \(J_ts+J_t^2(1/m+R^2/I)/2=-ms^2/7\)，故总损失为 \(K_{\rm before}-K_{\rm after}=m(1-e^2)u^2/2+ms^2/7\ge0\)。这同时检查了冲量方向及接触点切向约束。
